\section*{Problems}

\subsection*{Problem Statements}

% Remember
\begin{problema}
\label{prob:85eb17b}
A fitted quadratic model is $\hat{\texttt{mpg}}=64.20-0.7634\,\texttt{hp}+0.003657\,\texttt{hp}^2$. Calculate predicted \texttt{mpg} for cars with \texttt{hp}=60, 100, and 180, and identify the approximate horsepower at which the model predicts minimum \texttt{mpg}.
\end{problema}

% Understand
\begin{problema}
\label{prob:5afaae8}
Fitting two models to the same data (\texttt{hp} versus \texttt{mpg}) gives $R^2_{\text{linear}}=0.080$ for $\texttt{mpg}=c_0+c_1\,\texttt{hp}$ and $R^2_{\text{quadratic}}=0.994$ for $\texttt{mpg}=c_0+c_1\,\texttt{hp}+c_2\,\texttt{hp}^2$. Interpret this dramatic difference and explain what it says about the true functional form.
\end{problema}

% Apply
\begin{problema}
\label{prob:e757ebc}
With $n=5$ \texttt{hp}/\texttt{mpg} observations, $R^2_{\text{linear}}=0.080$ and $R^2_{\text{quadratic}}=0.994$. Perform the partial $F$ test ($H_0:c_2=0$, so the quadratic term is unnecessary) at $\alpha=0.05$ (critical value $F_{0.05,1,2}=18.51$), using $F_{\text{partial}}=\dfrac{(R^2_{\text{quadratic}}-R^2_{\text{linear}})/(2-1)}{(1-R^2_{\text{quadratic}})/(5-2-1)}$.
\end{problema}

% Analyze
\begin{problema}
\label{prob:c772cd1}
Prove that least-squares fitting of $Y=c_0+c_1X+c_2X^2+\varepsilon$ is formally identical to multiple linear regression if the auxiliary variable $Z=X^2$ is defined and $Y=c_0+c_1X+c_2Z+\varepsilon$ is fitted. Explain why the normal equation $\hat{\mathbf c}=(\mathbf X^T\mathbf X)^{-1}\mathbf X^T\mathbf Y$ applies without modification even though the model is nonlinear in the original $X$.
\end{problema}

% Evaluate
\begin{problema}
\label{prob:6954b16}
\textbf{Inherent multicollinearity in polynomial models:} in $Y=c_0+c_1X+c_2X^2+\varepsilon$, where $X$ is strictly positive and concentrated in a narrow range (for example $X\in[95,105]$), prove that the design columns $X$ and $X^2$ are highly correlated (correlation near 1), and evaluate why this inflates both predictors' VIF without invalidating model predictions, although it complicates isolated interpretation of $c_1$ and $c_2$.
\end{problema}

% Create
\begin{problema}
\label{prob:bc664de}
Design an original polynomial-regression problem (different from the \texttt{hp}-\texttt{mpg} relationship): a pricing economist models product units sold ($Y$) as a function of sale price ($X$, dollars), observing nonmonotone demand—relatively high sales at low prices, minimum sales at intermediate prices, and a rebound at very high prices because of exclusivity. The fitted model is $\hat Y=850-12X+0.09X^2$. Calculate predictions for $X=30$, 70, and 110, and identify the price at which minimum sales are predicted.
\end{problema}

\subsection*{Detailed Solutions}

% Remember
\begin{solproblema}[prob:85eb17b]
\begin{align}
\hat{\texttt{mpg}}(60)&=64.20-0.7634(60)+0.003657(60)^2\approx31.57,\\
\hat{\texttt{mpg}}(100)&\approx24.43,\\
\hat{\texttt{mpg}}(180)&\approx45.32.
\end{align}
The parabola's minimum is at $\texttt{hp}^*=-c_1/(2c_2)=0.7634/(2\times0.003657)\approx104.4$, so minimum \texttt{mpg} is predicted near \texttt{hp}=104.
\end{solproblema}

% Understand
\begin{solproblema}[prob:5afaae8]
The change from $R^2=0.080$ to $R^2=0.994$ is dramatic and shows that the \texttt{hp}-\texttt{mpg} relationship is \textbf{not linear}: the linear model explains only 8\% of variability, while the quadratic model explains 99.4\%. The underlying functional form is quadratic (or at least requires curvature), so conclusions based only on the linear model would be misleading.
\end{solproblema}

% Apply
\begin{solproblema}[prob:e757ebc]
\begin{align}
F_{\text{partial}}=\frac{(0.994-0.080)/(2-1)}{(1-0.994)/(5-2-1)}=\frac{0.914}{0.003}\approx304.7.
\end{align}
Since $304.7\gg18.51$, \textbf{reject $H_0$} decisively: the quadratic term $c_2\,\texttt{hp}^2$ is highly significant and necessary; the simple linear model is badly misspecified.
\end{solproblema}

% Analyze
\begin{solproblema}[prob:c772cd1]
Defining $Z=X^2$ rewrites the model exactly as $Y=c_0+c_1X+c_2Z+\varepsilon$, which is \emph{linear in the parameters} $c_0,c_1,c_2$. The design matrix includes a column of already computed $X^2$ values: $\mathbf X_{\text{design}}=[\mathbf1,X,X^2]$. Normal equations require linearity in the parameters, not linearity between $Y$ and the original predictor $X$, so they apply unchanged.
\end{solproblema}

% Evaluate
\begin{solproblema}[prob:6954b16]
For $X$ concentrated around a large central value $X_0=100$, write $X=X_0+\delta$ with $|\delta|\ll X_0$. Then $X^2=X_0^2+2X_0\delta+\delta^2\approx X_0^2+2X_0\delta$, so $X^2$ is approximately linear in $X$ over the narrow range and $r_{X,X^2}$ is close to 1. By $\text{VIF}_j=1/(1-R_j^2)$, the high correlation produces a large VIF for both predictors. This does not invalidate predictions, which depend on the stable combination $\hat c_1X+\hat c_2X^2$, but it compromises isolated interpretation of $\hat c_1$ and $\hat c_2$ because their standard errors inflate.
\end{solproblema}

% Create
\begin{solproblema}[prob:bc664de]
\begin{align}
\hat Y(30)&=850-12(30)+0.09(30)^2=571,\\
\hat Y(70)&=850-12(70)+0.09(70)^2=451,\\
\hat Y(110)&=850-12(110)+0.09(110)^2=619.
\end{align}
The parabola's minimum occurs at
\begin{align}
X^*=-\frac{c_1}{2c_2}=\frac{12}{2(0.09)}=\frac{12}{0.18}\approx66.67,
\end{align}
so the model predicts minimum sales near a price of \$66.67, consistent with the closest evaluated point, $\hat Y(70)=451$.
\end{solproblema}
